% =============================================================
% CHAPTER 5 : CONCLUSION
% =============================================================

\section{Conclusion}
\label{sec:conclusion}

% ---------------------------------------------------------
% 5.1 SUMMARY OF FINDINGS
% ---------------------------------------------------------

\subsection{Summary of Findings}
\label{sec:summary_of_findings}

This dissertation developed and applied a three-stage statistical
framework for health insurance cost prediction, combining deep neural
network regression, parametric Monte Carlo simulation and Shapley-value
attribution within a single empirical workflow.  The analysis was
conducted on a publicly available Kaggle benchmark dataset containing one
million policyholder records with eleven predictor variables and a
continuous monetary response.

The predictive stage demonstrated that a feed-forward deep neural network
with a funnel architecture, CELU activation and Nadam optimisation
achieves strong out-of-sample performance on the benchmark data, with
test-set $R^2 = 0.9957$, MAE $= \text{R}251$ and RMSE $= \text{R}290$.
A five-phase sequential hyperparameter comparison identified this
configuration from among 16 candidate settings.  A baseline comparison
with a Gamma generalised linear model (log-link) confirmed that the
\gls{dnn} substantially outperforms the actuarial-standard model class on
this dataset, reducing the test RMSE by $68.2\%$ and the test MAE by
$59.8\%$ relative to the \gls{glm}.  The near-identical training and test
performance metrics for both models ruled out overfitting, and the
data-leakage prevention protocols described in
Section~\ref{sec:data_leakage_prevention} ensured that all reported
metrics reflect genuine out-of-sample predictive ability.

The simulation stage generated $B = 10\,000$ synthetic policyholder
profiles from fitted parametric marginal distributions and propagated them
through the trained network to produce a model-implied predicted-cost
distribution centred at R16\,707 with moderate dispersion
(SD $=$ R4\,409) and a mildly right-skewed shape.  Convergence diagnostics
confirmed adequate Monte Carlo precision, and joint-distribution validation
(including an energy distance permutation test, $p = 0.730$) established
that the independent-marginal sampling does not introduce material
distributional distortion.  The residual-augmented sensitivity analysis
confirmed that these findings are robust to fitted residual perturbation.

The explanation stage computed gradient-based Shapley values for all
simulated profiles and summarised feature importance at both the global
and conditional levels.  At the global level, smoking status is the
dominant contributor to predicted cost ($I_g = 0.577$), followed by
coverage level ($0.401$), medical history ($0.368$) and family medical
history ($0.352$).  The conditional analysis at three upper-tail
thresholds ($\tau_{90}$, $\tau_{95}$, $\tau_{99}$) revealed a progressive
shift in the feature-importance ranking.  At $\tau_{99}$, family medical
history ($I_{g,0.99} = 0.622$), medical history ($0.608$) and coverage
level ($0.599$) overtake smoking status ($0.566$) as the most important
variables.  This reversal occurs because smoking status is near-universal
($99\%$) in the extreme tail and therefore loses its discriminating power;
the residual variation among the costliest profiles is instead explained by
coverage tier, medical history and family medical history.  A Shapley
attribution invariance check confirmed that these rankings are consistent
across training-data and Monte Carlo input regimes.

% ---------------------------------------------------------
% 5.2 CONTRIBUTIONS
% ---------------------------------------------------------

\subsection{Contributions}
\label{sec:contributions}

The dissertation makes the following contributions:

\begin{enumerate}
    \item \textbf{Integrated workflow.}  The study presents prediction,
          simulation and explanation as a single, reproducible pipeline
          rather than as isolated modelling exercises.  The integration
          allows the Monte Carlo distribution to define the regions of
          interest for the Shapley analysis, and the Shapley analysis to
          decompose the model behaviour within those regions.

    \item \textbf{Conditional Shapley analysis in the upper tail.}  The
          conditional feature-importance summaries at high-cost thresholds
          demonstrate that the variables driving extreme predicted costs
          differ from those that are most important on average.  This
          finding is relevant to actuarial practice, where understanding
          the drivers of high-cost claims is of particular interest.

    \item \textbf{Model-implied distribution under an explicit input law.}
          The Monte Carlo component constructs a model-implied
          predicted-cost distribution from a formally stated input
          distribution, rather than reporting predictions only at observed
          data points.  This connects the neural network prediction
          function to the simulation-based inference tradition within
          mathematical statistics.

    \item \textbf{Comprehensive validation.}  The dissertation includes
          joint-distribution diagnostics, convergence monitoring,
          residual-augmented sensitivity analysis and Shapley attribution
          invariance checks that go beyond standard predictive-performance
          reporting.

    \item \textbf{GLM benchmark contextualisation.}  The Gamma/log-link
          \gls{glm} comparison contextualises the \gls{dnn}'s performance
          within the actuarial modelling landscape and establishes that the
          predictive superiority reflects genuine model-class advantages
          rather than artefacts of the benchmark data.
\end{enumerate}

% ---------------------------------------------------------
% 5.3 STRENGTHS OF THE STUDY
% ---------------------------------------------------------

\subsection{Strengths of the Study}
\label{sec:strengths}

The study has the following methodological strengths:

\begin{enumerate}
    \item \textbf{Large benchmark dataset.}  The analysis is based on one
          million policyholder records, which provides a substantial sample
          for training, validation and testing.  The training partition
          alone contains $524\,288$ observations, reducing sampling
          variability in model fitting and Shapley-value estimation.

    \item \textbf{Systematic hyperparameter selection.}  The five-phase
          sequential comparison provides a transparent, reproducible method
          for selecting the network configuration.  Each phase is
          documented with validation-loss curves and selection criteria,
          allowing the choices to be scrutinised and replicated.

    \item \textbf{Rigorous validation pipeline.}  The data-leakage
          prevention protocols (Section~\ref{sec:data_leakage_prevention}),
          the joint-distribution input validation, the convergence
          diagnostics and the residual-augmented sensitivity analysis
          collectively establish that the Monte Carlo and Shapley results
          are not artefacts of data contamination, distributional
          mismatch or residual instability.

    \item \textbf{Multi-level attribution analysis.}  The Shapley-value
          analysis is conducted at both the global and conditional levels,
          with three upper-tail thresholds providing a graded view of how
          feature importance evolves into the extreme tail.  The attribution
          invariance check further validates the stability of the results
          across input regimes.

    \item \textbf{Transparent interpretive boundaries.}  All conclusions
          are explicitly qualified as internal to the fitted model and
          benchmark data.  No causal claims are made, and the modelling
          assumptions are stated formally in a consolidated table
          (Table~\ref{tab:modelling_assumptions_ch3}).
\end{enumerate}

% ---------------------------------------------------------
% 5.4 LIMITATIONS
% ---------------------------------------------------------

\subsection{Limitations}
\label{sec:limitations_conclusion}

The following limitations qualify the scope of the findings:

\begin{enumerate}
    \item \textbf{Benchmark data.}  All results are conditional on the
          Kaggle \textit{Insurance Data for Machine Learning} benchmark
          dataset, which has a strong simulated structure.  The unusually
          high predictive performance ($R^2 = 0.9957$) reflects this
          structure and should not be extrapolated to real-world insurance
          portfolios, where data are noisier and relationships between
          covariates and costs are more complex.

    \item \textbf{Limited model comparison.}  The study compares the
          \gls{dnn} against a single classical baseline (Gamma \gls{glm}
          with log-link).  Tree-based models such as random forests and
          gradient-boosted trees, ensemble methods and Bayesian
          alternatives were not evaluated.  A broader model-selection study
          would strengthen the case for the \gls{dnn} as the preferred
          prediction surface.

    \item \textbf{Single fitted network.}  The Monte Carlo and Shapley
          analyses are conditional on a single trained network with a
          fixed random seed.  Variability across random initialisations
          was not quantified.

    \item \textbf{Independent marginals.}  The Monte Carlo input
          distribution is constructed as a product of independent
          marginals.  While the joint-distribution validation confirmed
          that the training-data dependence structure is weak, the
          independence assumption broadens the simulated input space beyond
          what would be observed in practice.

    \item \textbf{Approximate Shapley values.}  The reported attributions
          are gradient-based SHAP approximations, not exact game-theoretic
          Shapley values.  They are conditional on the fitted network, the
          background sample of $m_0 = 200$ observations and the explainer
          settings.

    \item \textbf{Small $\tau_{99}$ sample.}  The conditional Shapley
          summaries at the 99th percentile threshold rely on approximately
          100 simulated profiles.  The ranking reversal observed at
          $\tau_{99}$ is consistent with the monotonic trends at
          $\tau_{90}$ and $\tau_{95}$, but its precision could be improved
          with larger simulation sizes or bootstrap confidence intervals.

    \item \textbf{No causal inference.}  The Shapley decomposition
          identifies variables to which the fitted model assigns large
          contributions.  It does not establish causal relationships
          between input variables and healthcare expenditure.
\end{enumerate}

% ---------------------------------------------------------
% 5.5 RECOMMENDATIONS FOR FUTURE RESEARCH
% ---------------------------------------------------------

\subsection{Recommendations for Future Research}
\label{sec:future_research}

Several directions could extend the framework developed in this
dissertation:

\begin{enumerate}
    \item \textbf{Broader model comparison.}  Future work could extend
          the baseline comparison to include tree-based ensembles (random
          forests, XGBoost, LightGBM) and Bayesian neural networks.  Such
          comparisons would clarify whether the \gls{dnn}'s predictive
          advantage over the \gls{glm} extends to other flexible model
          classes, and whether the Shapley attribution patterns are
          model-dependent or model-invariant.

    \item \textbf{Copula-based Monte Carlo inputs.}  Replacing the
          product-marginal input distribution with a copula-based joint
          distribution would allow the simulation to reproduce the
          dependence structure of the observed data more faithfully.  This
          would be particularly relevant for datasets where the input
          variables exhibit meaningful correlations.

    \item \textbf{Robustness of extreme-tail attributions.}  The
          $\tau_{99}$ ranking reversal could be strengthened by increasing
          the simulation size to $B = 50\,000$ or $100\,000$, by
          constructing bootstrap confidence intervals for the conditional
          mean absolute Shapley values, or by introducing intermediate
          thresholds (e.g.\ $\tau_{97.5}$, $\tau_{98}$) to trace the
          transition in the importance ranking more finely.

    \item \textbf{Multi-seed stability analysis.}  Training the
          \gls{dnn} under multiple random initialisations and reporting
          the distribution of Shapley-value summaries across seeds would
          quantify the sensitivity of the attribution results to the
          stochastic training process.

    \item \textbf{Real-world actuarial data.}  Applying the framework to
          proprietary or regulatory insurance datasets with genuine
          underwriting heterogeneity would test whether the conditional
          Shapley findings generalise beyond the synthetic benchmark
          setting.

    \item \textbf{Bayesian and conformal prediction intervals.}  Coupling
          the Monte Carlo simulation with Bayesian posterior predictive
          distributions or conformal prediction intervals would add formal
          uncertainty quantification to the predicted-cost distribution.
\end{enumerate}

% ---------------------------------------------------------
% 5.6 CONCLUDING REMARK
% ---------------------------------------------------------

\subsection{Concluding Remark}
\label{sec:concluding_remark}

This dissertation has demonstrated that combining deep neural network
prediction, Monte Carlo simulation and Shapley-value attribution within a
single workflow produces richer insight into a fitted model's behaviour
than any one component in isolation.  The predicted-cost distribution
characterises the range and shape of the model's output landscape under a
specified input law, and the conditional Shapley analysis reveals that the
factors driving extreme predicted costs are not the same as those that
dominate on average.  These findings contribute to the intersection of
statistical learning, simulation-based inference and explainable modelling,
and provide a methodological template that can be adapted to other
regression problems where distributional analysis and feature attribution
in specific regions of the output space are of interest.
